\documentclass[12pt,a4paper]{article}
\usepackage[margin=2.4cm]{geometry}
\usepackage{amsmath,amssymb,bm,mathtools}
\usepackage{fontspec}
\setmainfont{TeX Gyre Termes}
\setsansfont{TeX Gyre Heros}
\setmonofont{Latin Modern Mono}[Scale=0.9]
\usepackage{graphicx,float,caption,subcaption,booktabs,longtable,array,multirow,tabularx}
\usepackage{xcolor}
\usepackage{setspace}
\usepackage{lineno}
\usepackage{microtype}
\usepackage{xurl}
\usepackage{siunitx}
\usepackage{enumitem}
\usepackage{fancyhdr}
\usepackage{titlesec}
\titleformat{\paragraph}[runin]{\normalfont\itshape}{}{0pt}{}
\titlespacing*{\paragraph}{0pt}{1.1ex plus .2ex}{0.7em}
\usepackage[authoryear,round,semicolon,sort&compress]{natbib}
\definecolor{lk}{RGB}{0,70,140}
\usepackage[colorlinks=true,linkcolor=lk,citecolor=lk,urlcolor=lk,breaklinks=true,pdfborder={0 0 0}]{hyperref}
\usepackage[nameinlink,capitalise]{cleveref}
\sisetup{detect-all,separate-uncertainty=false}
\setstretch{1.5}
\setlength{\parskip}{3pt}
\captionsetup{font=small,labelfont=bf,justification=centering,singlelinecheck=true}
\titleformat{\section}{\large\bfseries}{\thesection.}{0.5em}{}
\titleformat{\subsection}{\normalsize\bfseries}{\thesubsection}{0.5em}{}
\titleformat{\subsubsection}{\normalsize\itshape}{\thesubsubsection}{0.5em}{}
\pagestyle{fancy}\fancyhf{}\fancyfoot[C]{\small\thepage}\renewcommand{\headrulewidth}{0pt}
\setlength{\headheight}{14pt}
\bibliographystyle{plainnathl}
\newcommand{\doihl}[1]{doi: \href{https://doi.org/#1}{\nolinkurl{#1}}}
\newcommand{\dd}{\mathrm{d}}
\newcommand{\ET}{\mathit{ET}}
\newcommand{\Ent}{\mathit{ET}_{\mathrm{nt}}}
\newcommand{\Ks}{K_{\mathrm{s}}}
\newcommand{\Se}{S_{\mathrm{e}}}
\newcommand{\TAW}{\mathit{TAW}}
\newcommand{\repourl}{https://github.com/halilunaziru73-creator/Coupled-Water-Fluxes-Pore-Scale-Dynamics-Precision-Irrigation-Mulching-Feedback-Loops}
\newcommand{\repo}{\href{\repourl}{repository}}
\newcommand{\dataurl}{\repourl/tree/main/Data}
\newcommand{\codeurl}{\repourl/tree/main/Manuscripts}
\newcommand{\softurl}{\repourl/tree/main/Software}
\newcommand{\datalink}{\href{\dataurl}{Data}}
\newcommand{\codelink}{\href{\codeurl}{Code}}
\newcommand{\softlink}{\href{\softurl}{Software}}
\newcommand{\manuallink}{\href{\manualurl}{Manual}}
\newcommand{\downloadlink}{\href{\exeurl}{SWAFL.exe}}
\newcommand{\figdir}{../figures}
\newcommand{\manualurl}{https://github.com/halilunaziru73-creator/Coupled-Water-Fluxes-Pore-Scale-Dynamics-Precision-Irrigation-Mulching-Feedback-Loops/raw/main/Software/v1.1.0/SWAFL_Desktop/docs/manual_v1.1.0/SWAFL_Manual_v1.1.0.pdf}
\newcommand{\exeurl}{https://github.com/halilunaziru73-creator/Coupled-Water-Fluxes-Pore-Scale-Dynamics-Precision-Irrigation-Mulching-Feedback-Loops/releases/latest/download/SWAFL.exe}
\newcommand{\authorblock}{%
\begin{center}
{\large\textbf{Naziru Halilu}}$^{1,2,*}$\\[3pt]
{\small $^{1}$ Agricultural University of Athens, Athens, Greece\\
$^{2}$ Public University of Navarre, Pamplona, Spain\\
ORCID: \href{https://orcid.org/0009-0006-2324-9599}{0009-0006-2324-9599}\\
$^{*}$ Correspondence: \href{mailto:halilunaziru73@gmail.com}{halilunaziru73@gmail.com}}
\end{center}}
\newcommand{\csurl}{https://github.com/halilunaziru73-creator/Coupled-Water-Fluxes-Pore-Scale-Dynamics-Precision-Irrigation-Mulching-Feedback-Loops/tree/main/Data/Comprehensive_Summary}
\newcommand{\cslink}{\href{\csurl}{\nolinkurl{Data/Comprehensive_Summary}}}
\newcommand{\csavail}{\par\smallskip\textbf{Full simulation outputs.} Every value of every log of every run (the annual water and nitrogen balances, the harvests and the crop state, year by year), one compressed table per experiment (\texttt{E00\_Baselines} to \texttt{E16\_Edge\_Cases}), the one-row-per-run summary of all 10\,190 runs (\texttt{00\_summary\_all\_simulations}), the per-experiment summary (\texttt{00\_SUMMARY\_ALL\_EXPERIMENTS}) and the script that made them (\texttt{summarize\_results3.py}) are in the folder \cslink{} of the repository (\url{\csurl}); their 20-year means reproduce the run table within 2\,\%.}
% Data and code availability of the two analysis papers (P1, P2)
\newcommand{\dataavail}{%
\textbf{Data.} The complete Daisy ensemble outputs (all 10\,190 runs, raw and derived) are archived on Zenodo, \doihl{10.5281/zenodo.23043626}. The raw run table used here, and a linked list of every data file used, are in the \datalink{} folder of the project repository (\url{\dataurl}).\csavail
}
% Software, data and code availability of the software paper (P3)
\newcommand{\softavail}{%
\textbf{Software.} SWAFL Desktop (version 1.1.0) is open source; the \softlink{} source is in the project repository (\url{\softurl}) and the program for Windows is \downloadlink{} (\url{\exeurl}).
\par\smallskip\textbf{Manual.} The \manuallink{} (\url{\manualurl}) documents every control, view and equation and labels the provenance of all 134 numbered equations; it can be reviewed or downloaded as a PDF.
\par\smallskip\textbf{Data.} The complete Daisy ensemble outputs (all 10\,190 runs, raw and derived) are archived on Zenodo, \doihl{10.5281/zenodo.23043626}. The raw run table, and a linked list of every data file used, are in the \datalink{} folder of the repository (\url{\dataurl}).\csavail
}

\usepackage{pdflscape}
\setlength{\emergencystretch}{3em}

% A shared availability statement
\newcommand{\suiteavail}{%
\textbf{Software.} SWAFL Desktop (version 1.1.0) is open source; the \softlink{} source is in the project repository (\url{\softurl}) and the program for Windows is \downloadlink{} (\url{\exeurl}). The whole workbench also runs in a browser, with nothing to install (link in the repository README).
\par\smallskip\textbf{Manual.} The \manuallink{} (\url{\manualurl}) documents every control, screen and equation of the release, with labelled screenshots.
\par\smallskip\textbf{Data.} The complete Daisy ensemble outputs (all 10\,190 runs, raw and derived) are archived on Zenodo, \doihl{10.5281/zenodo.23043626}; the run table and a summary of every experiment are in the \datalink{} folder (\url{\dataurl}).\csavail
}

\renewcommand{\thesection}{S\arabic{section}}
\renewcommand{\thetable}{S\arabic{table}}
\hypersetup{pdftitle={Supplementary Material 5 The Open Software Stack Behind SWAFL Desktop},pdfauthor={Naziru Halilu}}
\usepackage{tikz}
\usepackage[most]{tcolorbox}
\newcommand*\circled[1]{\tikz[baseline=(c.base)]\node[circle,draw=red!70!black,text=red!70!black,line width=0.6pt,inner sep=0.7pt,minimum size=1.15em,font=\scriptsize\bfseries\sffamily](c){#1};}
\newtcolorbox{problem}[1][]{enhanced,breakable,colback=orange!6,colframe=orange!70!black,fonttitle=\bfseries\sffamily,title=The problem,#1}
\newtcolorbox{derivation}[1][]{enhanced,breakable,colback=white,colframe=black!55,fonttitle=\bfseries\sffamily,title=The derivation,#1}
\newtcolorbox{keyresult}[1][]{enhanced,breakable,colback=green!5,colframe=green!40!black,fonttitle=\bfseries\sffamily,title=The result,#1}
\newtcolorbox{pins}[1][]{enhanced,breakable,colback=blue!3,colframe=blue!50!black,fonttitle=\bfseries\sffamily,title=Pinpoints: what each circled mark shows and why,#1}
\newcommand{\pin}[2]{\par\noindent\circled{#1}\ #2\par\smallskip}

\begin{document}
\begin{center}{\large\bfseries Supplementary Material 5\par}\vspace{4pt}
{\Large\bfseries The Open Software Stack Behind SWAFL Desktop\par}\vspace{4pt}
{\normalsize Languages and libraries and code and techniques and tools used to build the software and its documents\par}\end{center}
\authorblock
\tableofcontents
\clearpage

\section{Purpose and scope}
This supplement stands alone. It belongs to the four articles of the series but is not tied to one of them: it records what SWAFL Desktop 1.1.0 is made of, in which languages it is written, which libraries and programs it carries, which techniques its parts use, and which tools were used to build, test and document it. Every figure below was counted on the source of version 1.1.0 or read from its \texttt{package.json}; where a number is approximate the text says so. Nothing in the stack is proprietary and nothing needs an account: all parts are open source or are the model and data that the series publishes.

\section{Read this first: the words explained}
Software is made like a meal in a restaurant kitchen. The words below use that picture; the rest of this supplement uses the same words with the same meaning.
{\small
\begin{longtable}{@{}>{\raggedright\arraybackslash}p{2.9cm}>{\raggedright\arraybackslash}p{4.9cm}>{\raggedright\arraybackslash}p{5.2cm}@{}}
\caption{The words of this supplement in plain language.}\label{tab:words}\\
\toprule Word & What it means & In the kitchen\\\midrule\endfirsthead
\toprule Word & What it means & In the kitchen\\\midrule\endhead
Language & The way instructions are written for a computer (Python and JavaScript are languages). & The language of the recipe.\\
Code & The instructions themselves, written in a language. & The written recipe.\\
Runtime & A program that reads the code and carries it out. Electron and Node.js and Python are runtimes. & The cook who can read that recipe.\\
Library & Code written by someone else that your code calls for one job (draw a map or read an Excel file). & A ready-made sauce or a stock you buy instead of making.\\
Package & A library or tool wrapped so that one command installs it (\texttt{pip install} or \texttt{npm install}). & The delivery box in which the sauce arrives.\\
Dependency & A package that another package needs in order to work. & The ingredients that the sauce itself is made from.\\
Program or tool & A finished application that you start (ffmpeg and pandoc and Daisy). & A kitchen machine such as an oven: you use it but you do not change it.\\
Framework & A library that also decides how your whole program must be arranged. SWAFL uses none. & A franchise kitchen that dictates the whole menu.\\
Compiler & A program that turns code into a file the computer runs directly. Daisy was compiled from C++. & Baking the dough into bread: you eat the bread, not the dough.\\
Version & A number that says which release of a package is meant (for example 2.5.3). & The year of the vintage.\\
Pinned & A version that is written down so that the same one is installed every time. & Always buying the same batch.\\
Licence & The permission that the author gives to use and copy the code. MIT and Apache and BSD allow free use. & The label that says you may cook with it.\\
Build & Putting the parts together into the finished program or document. & Plating the dish.\\
Bridge & The small piece of code that lets two parts talk to each other. & The waiter between the kitchen and the table.\\
Service (online) & A website of another organisation that answers questions (weather or satellite pictures). & Ordering an ingredient by phone.\\
CDN & A web server that hands out libraries to many websites. SWAFL does not use one at run time. & A supplier who delivers every day; here the shelf is already stocked.\\
\bottomrule\end{longtable}}

\section{How the program is made, step by step in plain words}
\begin{enumerate}[leftmargin=1.5em]
\item \textbf{The author writes code} in several languages: JavaScript for what the user sees and does, Python for the scripts that prepare data and figures, and \LaTeX\ and Markdown for the articles and the manual (Table~\ref{tab:lang}).
\item \textbf{The code borrows libraries} for jobs that are not worth writing again: a map (Leaflet), satellite file reading (GeoTIFF.js), Excel reading (SheetJS) (Table~\ref{tab:libs}). Python scripts borrow packages for mathematics and figures (Table~\ref{tab:py}).
\item \textbf{The model is Daisy.} Daisy is an existing, tested scientific program written by others in C++. SWAFL does not copy or change it; it writes the question into a file, starts Daisy, and reads the answer back from the files Daisy writes.
\item \textbf{Electron wraps everything} into one window, like putting the kitchen and the dining room in one building. It carries its own browser and Node.js, so the user does not need to install either.
\item \textbf{The build tool packs one file.} \texttt{electron-builder} packs the code, the libraries, Daisy and the manual into \texttt{SWAFL.exe} (Table~\ref{tab:build}).
\item \textbf{The documents are built the same way:} the manual from Markdown by pandoc and XeLaTeX, and the articles from \LaTeX, with figures made by the Python scripts.
\item \textbf{Tests run the real program} in a hidden browser (Playwright) and check that no page error appears and that the numbers agree with independent references.
\end{enumerate}

\section{Who needs what}
\begin{table}[H]\centering\small
\caption{What each kind of reader has to install.}\label{tab:who}
\begin{tabular}{@{}>{\raggedright\arraybackslash}p{3.4cm}>{\raggedright\arraybackslash}p{9.6cm}@{}}
\toprule Reader & What to do\\\midrule
A user of the program & Nothing to install. Double click \texttt{SWAFL.exe}. Everything it needs is inside the file, and the online services are optional.\\
A reader who only wants to look at the code & Open the folders in a web browser or a text editor. No installation.\\
A researcher who repeats the numbers & Install Python 3.13 and the packages of the pip section, then run \texttt{run\_all.sh}.\\
A developer who rebuilds the program & Install Node.js 22 and run \texttt{npm install} and \texttt{npm run get-app}.\\
A person who rebuilds the manual and articles & Install pandoc, TeX Live (XeLaTeX and BibTeX) and the fonts of the programs section.\\
\bottomrule\end{tabular}\end{table}

\section{The software at a glance}
SWAFL Desktop is one program with four layers, from the outside in.
\begin{enumerate}[nosep,leftmargin=1.5em]
\item \textbf{The desktop shell.} An Electron application (\texttt{main.js}, \texttt{preload.js}, \texttt{daisy-runner.js}, \texttt{serve.js}) that opens the window, saves and opens files, starts Daisy and opens the manual.
\item \textbf{The workbench.} About 90 JavaScript files in \texttt{app/js} (layers, properties, toolbox, dashboard, field on the map, sketch, the Daisy tabs, scientific figures, reports).
\item \textbf{The simulation engine.} A single page (\texttt{engine.html}, 1.5\,MB, 7\,225 lines) with the scripts of \texttt{app/engine} (crops, weather, the layers of the seasons, deterministic phenotyping). It runs in an inline frame and is joined to the workbench by a bridge.
\item \textbf{The Daisy model.} Daisy 7.1.14 for Windows with its library, packed unchanged beside the program and started as a separate process.
\end{enumerate}
The same workbench and engine run in a web browser from a small local server (\texttt{npm run web}, or the \texttt{START\_SWAFL\_WEB} files) and online from the project page, so one code base serves the desktop, the browser and the web.

\section{Code map of the repository}
Code is kept in four folders of the repository and nothing is hidden outside them. The counts are of the files tracked in the repository on the day of this supplement (the Daisy program folder and the vendored libraries are not counted).
\begin{table}[H]\centering\small
\caption{Where the code is, by folder.}\label{tab:codemap}
\begin{tabular}{@{}>{\raggedright\arraybackslash}p{4.3cm}>{\raggedright\arraybackslash}p{1.3cm}>{\raggedright\arraybackslash}p{1.3cm}>{\raggedright\arraybackslash}p{6.3cm}@{}}
\toprule Folder & Python files & JS files & What the code does and why it is there\\\midrule
\texttt{Software/\allowbreak v1.1.0/\allowbreak SWAFL\_\allowbreak Desktop} & 20 & 90 & the program (shell and workbench and engine) and the tools that build its data and manual (14\,100 lines of JavaScript and 2\,500 of Python)\\
\texttt{Software/\allowbreak Old\_\allowbreak Version\_\allowbreak 1.0} & 16 & 12 & the archived version 1.0 with its own build files (kept for reference)\\
\texttt{Manuscripts} & 75 & 27 & figures and tables and checks of the articles (\texttt{reproducibility}) and the graphical abstracts; 103 \LaTeX\ files; scripts for the screenshots and checks in Node.js with Playwright\\
\texttt{Data} & 4 & 0 & the workbook of the data and the calculation sheets\\
\texttt{Results} & 13 & 0 & the 50 charts, the interactive charts, the videos, the 3D project and the Excel analysis (\texttt{08\_Scripts})\\
\bottomrule\end{tabular}\end{table}

\section{Languages}
A language is only the way the instructions are written; the column on the right says what each one is used for in plain terms.
\begin{table}[H]\centering\small
\caption{Languages and formats of the 1.1.0 source (the Daisy program and the manual PDF are not counted).}\label{tab:lang}
\begin{tabular}{@{}>{\raggedright\arraybackslash}p{3.3cm}>{\raggedright\arraybackslash}p{2.0cm}>{\raggedright\arraybackslash}p{2.1cm}>{\raggedright\arraybackslash}p{6.6cm}@{}}
\toprule Language or format & Files & Lines (about) & Role and why it is used in the build\\\midrule
JavaScript (ECMAScript) & 90 & 14\,100 & workbench and engine scripts and desktop shell; every equation of the layers and designs. Why: it runs in the browser and in Electron without a compiler, so one code base serves the desktop and the web.\\
HTML and inline script & 2 & 7\,200 & the window (\texttt{index.html}) and the simulation engine (\texttt{engine.html}). Why: the engine is one page that can be opened alone and tested alone.\\
CSS & 4 & 630 & the classic grey and navy look and the dashboard. Why: one style sheet keeps the same classic look in every window.\\
Python 3.13 & 20 & 2\,500 & build tools of the data and of the manual (indexes and annotated screenshots and the equation table). Why: quick scripts for data, figures and file building that need no packaging.\\
JSON & 4 & 20\,500 & libraries and indexes read by the program (for example the manual index). Why: readable data that JavaScript reads directly and that a person can edit.\\
Markdown & 47 & 31\,000 & the chapters of the manual. Why: the text of the manual stays plain text under version control and builds to one PDF.\\
\LaTeX\ & 103 & -- & the four articles and their supplements (XeLaTeX with BibTeX). Why: exact equations and numbered references, and the same source gives the article and its supplement.\\
Batch and shell scripts & 7 & 80 & start and build files for Windows and macOS and Linux. Why: a double click starts the program without the user opening a terminal.\\
C++ (Daisy) & -- & -- & the source language of the Daisy model; used here only as the compiled program of Daisy 7.1.14. Why: it is the model whose runs the suite reads, so it is not rewritten.\\
SVG and CSV and XLSX & 4 and 24 and 2 & -- & vector figures and plain tables and Excel workbooks. Why: open formats that any program can read.\\
Daisy log formats (\texttt{.dlf} and \texttt{.dwf}) & 29 & -- & the daily logs and weather files that Daisy writes and reads. Why: the bridge reads them exactly as Daisy writes them.\\
Daisy set-up language (\texttt{.dai}) & library & -- & the model description that Daisy reads; written by the workbench and by the user. Why: it is the only way to give Daisy a scenario, so the workbench writes it for the user.\\
\bottomrule\end{tabular}\end{table}
The model of Daisy itself is written in C++ and is used here as a compiled program; SWAFL does not change its source. No framework (React or Vue or Angular) and no bundler is used: the pages are plain JavaScript loaded in order, so that the program can be read and changed with a text editor.

\section{Runtime stack}
\begin{table}[H]\centering\small
\caption{What runs when SWAFL.exe is opened.}\label{tab:runtime}
\begin{tabular}{@{}>{\raggedright\arraybackslash}p{3.6cm}>{\raggedright\arraybackslash}p{2.8cm}>{\raggedright\arraybackslash}p{6.6cm}@{}}
\toprule Part & Version & Role and why it is used in the build\\\midrule
Electron & 31 (range \texttt{\^{}31.0.0}) & window and Chromium and Node.js in one package; the portable build is made with it. Why: it packs a browser and Node.js into one program so that the web code runs as a desktop application without installation.\\
Chromium and Node.js & those of Electron 31 & draw the canvases and run the JavaScript; Node gives file and process access to the shell only. Why: Chromium draws the canvases and Node.js lets the shell read and write files and start Daisy.\\
Daisy & 7.1.14 & soil and plant and atmosphere process model (University of Copenhagen); its Windows folder carries the program, a Python 3.13 runtime for its Python functions and the C++ libraries it needs. Why: the process model whose runs the decision layers show; it is used as built, not rewritten.\\
\bottomrule\end{tabular}\end{table}
The shell and the page talk through a small bridge (\texttt{preload.js}, \texttt{contextBridge}) with about 20 named handlers (save a file, open a file, open the manual, find Daisy, run a Daisy job, read a log, list the runs, fetch an address, and a few window commands). The page itself has no direct access to the disk.

\section{Libraries carried inside the program}
\begin{table}[H]\centering\small
\caption{Third party code that is copied into the program (all under open licences; the licence text sits beside each file).}\label{tab:libs}
\begin{tabular}{@{}>{\raggedright\arraybackslash}p{3.0cm}>{\raggedright\arraybackslash}p{1.9cm}>{\raggedright\arraybackslash}p{2.6cm}>{\raggedright\arraybackslash}p{5.5cm}@{}}
\toprule Library & Version & Licence & Role and why it is used in the build\\\midrule
Leaflet & 1.9.4 & BSD 2-Clause & the map of the field (satellite and street layers; outlining the field). Why: a small map library that works offline with local code and gives the field outline tool.\\
GeoTIFF.js & bundled & MIT & reading the Sentinel-2 bands of the field from cloud-optimised GeoTIFF files in the browser. Why: satellite bands can be read in the page by range request, so no server of ours is needed.\\
SheetJS (xlsx) & 0.18.5 & Apache 2.0 & importing weather from Excel and CSV and writing result tables. Why: weather and results travel in Excel files, which users already have.\\
Daisy 7.1.14 with Boost and SuiteSparse and the GCC runtime & as shipped & GPL and LGPL families (see the Daisy folder) & the process model and its numerical libraries; started as a separate program, not linked into SWAFL. Why: these are what Daisy itself needs to run, so they are shipped with it.\\
\bottomrule\end{tabular}\end{table}
Everything else, including the figures, the report writer, the equations of the layers and the pipe hydraulics, is code written for this program.

\section{Libraries of the whole repository, by language}
\subsection{Python}
All third party Python packages imported anywhere in the repository (\texttt{Software}, \texttt{Manuscripts}, \texttt{Data}, \texttt{Results}) are listed. The versions of the first group are those recorded in \texttt{Manuscripts/reproducibility/requirements.txt} (Python 3.13.16 on Linux x86-64); for the others no version is recorded in the repository and none is claimed. The standard library (os and sys and re and json and codecs and subprocess and others) is used throughout and is not listed.
\begin{table}[H]\centering\small
\caption{Python packages and what each is used for.}\label{tab:py}
\begin{tabular}{@{}>{\raggedright\arraybackslash}p{2.7cm}>{\raggedright\arraybackslash}p{1.7cm}>{\raggedright\arraybackslash}p{8.6cm}@{}}
\toprule Package & Version & Role and why it is used in the build\\\midrule
numpy & 2.5.3 & arrays and the numerics of every analysis script. Why: fast arrays so that the analysis of 10\,190 runs runs in seconds.\\
scipy & 1.18.1 & statistics (Spearman correlation and tests) and curve fitting and root finding (\texttt{brentq}, \texttt{curve\_fit}, \texttt{minimize}) and \texttt{solve\_ivp} and sparse matrices (the checks of the closed forms). Why: tested solvers and fitting routines instead of hand-written ones, which makes the checks independent.\\
pandas & 3.0.5 & reading the ensemble of 10\,190 runs and summaries. Why: tables of runs are grouped and joined by name, which removes indexing mistakes.\\
scikit-learn & 1.9.1 & gradient boosting and random forest regression with grouped cross validation (\texttt{GroupKFold}), partial dependence and PCA on the ensemble. Why: standard learners and cross validation give results that others can repeat.\\
statsmodels & 0.15.0 & ordinary least squares models and their confidence intervals. Why: least squares with full statistical output for the effects reported in the articles.\\
matplotlib & 3.11.2 & every figure; \texttt{mpl\_toolkits.mplot3d} for the 3D project; \texttt{animation} with the FFMpegWriter for the videos. Why: one library makes the print figures and the videos with a shared style.\\
Pillow (PIL) & 12.3.0 & screenshots and figure assembly and annotation. Why: puts screenshots and panels together without a graphics program.\\
openpyxl & not pinned & writing the Excel workbooks of the data and analysis. Why: writes the Excel files of the data with formulas and formats.\\
plotly & not pinned & the interactive charts and the 3D project saved as web pages. Why: charts that can be zoomed and rotated in a web page without a server.\\
pypdf & not pinned & preparing A4 versions of the articles. Why: sets the page size and joins pages of the PDFs.\\
wordfreq & not pinned & word frequencies (\texttt{zipf\_frequency}) for a check of the wording in the analysis folder. Why: a check that the wording is plain and common.\\
playwright (Python) & not pinned & screenshots of the program for the manual. Why: takes the screenshots of the real program, so the manual shows what the user sees.\\
\bottomrule\end{tabular}\end{table}
Small modules written for the repository (\texttt{common.py}, \texttt{load.py}, \texttt{figstyle.py}, \texttt{p\_style.py}, \texttt{pm\_style.py}, \texttt{annot.py}, \texttt{spec\_atlas.py}, \texttt{textjs.py} and the like) hold the shared figure style and the loaders, so that every figure of the series has the same look.

\subsection{JavaScript and HTML}
\begin{table}[H]\centering\small
\caption{JavaScript and HTML libraries and runtimes, end to end.}\label{tab:js}
\begin{tabular}{@{}>{\raggedright\arraybackslash}p{3.1cm}>{\raggedright\arraybackslash}p{1.9cm}>{\raggedright\arraybackslash}p{7.9cm}@{}}
\toprule Library or runtime & Version & Role and why it is used in the build\\\midrule
Electron & 31 & shell of the desktop program (\texttt{require('electron')}: app and window and ipcMain and contextBridge and dialogs). Why: it packs a browser and Node.js into one program so that the web code runs as a desktop application without installation.\\
electron-builder & 24.13.3 & development only: packs \texttt{SWAFL.exe} (portable) and the installer. Why: the standard way to make a Windows portable file and installer from the same folder.\\
Node.js built-in modules & of Node 22 and Electron 31 & \texttt{fs} and \texttt{path} and \texttt{child\_process} and \texttt{os} and \texttt{zlib} and \texttt{http} and \texttt{vm} and \texttt{url} and \texttt{crypto}: the local server, the Daisy runner, the build and test scripts. Why: no extra package is needed for files and a local server, so the program has no hidden dependencies.\\
Leaflet & 1.9.4 & the map of the field in the program. Why: a small map library that works with local code and gives the field outline tool.\\
GeoTIFF.js & bundled & Sentinel-2 bands read in the browser. Why: satellite bands can be read in the page by range request, so no server of ours is needed.\\
SheetJS (xlsx) & 0.18.5 & Excel and CSV import and export. Why: weather and results travel in Excel files, which users already have.\\
Playwright for Node & not pinned & test and screenshot scripts that drive the real program in headless Chromium. Why: drives the real program as a user would, so a page error is caught before release.\\
Plotly.js & 4.1.1 & the interactive charts and the 3D project in \texttt{Results} (the library file is stored beside the pages so that they work offline). Why: the engine of the interactive charts, stored locally so they work offline.\\
\bottomrule\end{tabular}\end{table}
Besides these, the program uses only what the browser provides: the 2D canvas, \texttt{Blob} and download links, \texttt{DecompressionStream} for compressed data, \texttt{localStorage} for small settings, \texttt{requestAnimationFrame} for the playing of a season, and \texttt{fetch} for the optional online services. There is no CDN call at run time: the libraries above are stored in the program folders.

\subsection{Programs called from the code}
\begin{table}[H]\centering\small
\caption{Command-line programs that scripts call.}\label{tab:progs}
\begin{tabular}{@{}>{\raggedright\arraybackslash}p{3.3cm}>{\raggedright\arraybackslash}p{9.7cm}@{}}
\toprule Program & Role and why it is used in the build\\\midrule
Daisy 7.1.14 (\texttt{daisy-bin.exe}) & started by \texttt{daisy-runner.js}; reads a set-up file and writes the logs. Why: it is the model itself, run as a separate process so a failed run cannot stop the program.\\
ffmpeg (libx264) & through matplotlib in \texttt{videos.py}: the videos of \texttt{Results/05\_Videos}. Why: matplotlib needs it to write the MP4 files.\\
pandoc & the manual: Markdown chapters to one PDF with the XeLaTeX engine. Why: turns the Markdown chapters into the PDF of the manual in one step.\\
XeLaTeX and BibTeX (TeX Live) & the manual and the articles and the supplements. Why: XeLaTeX uses system fonts and Unicode, and BibTeX keeps the references consistent.\\
Tesseract and poppler (\texttt{pdftoppm}) & checking and annotating figures in \texttt{annotate\_figs.py}. Why: reads text in figures and turns PDF pages into pictures for the checks.\\
Chromium (through Playwright) & screenshots and the tests of the program. Why: the same engine as the program, so the tests see the same pictures as the user.\\
\bottomrule\end{tabular}\end{table}

\subsection{\LaTeX\ packages and fonts}
The articles and supplements load these packages (from \texttt{pre.tex} and \texttt{supp\_style.tex}): amsmath and amssymb and array and bm and booktabs and caption and cleveref and enumitem and fancyhdr and float and fontspec and geometry and graphicx and hyperref and lineno and longtable and mathtools and microtype and multirow and natbib and pdflscape and setspace and siunitx and subcaption and tabularx and tcolorbox and tikz and titlesec and xcolor and xurl. The fonts are TeX Gyre Termes for the text and TeX Gyre Heros for sans serif and Latin Modern Mono for code. The manual uses DejaVu Serif and DejaVu Sans Mono.

\section{Installing everything: the exact commands}
These commands install every package of the tables above under the names that the package indexes use (note that scikit-learn is imported as \texttt{sklearn} and Pillow as \texttt{PIL}, and PyMuPDF as \texttt{pymupdf}). Python 3.13 and Node.js 22 are assumed. A user of the finished program needs none of this: \texttt{SWAFL.exe} carries what it needs.

\subsection{Python packages (pip)}
\begin{verbatim}
python3 -m pip install -r Manuscripts/reproducibility/requirements_all.txt
python3 -m playwright install chromium
\end{verbatim}
which is the same as
\begin{verbatim}
python3 -m pip install numpy==2.5.3 scipy==1.18.1 pandas==3.0.5 \
    scikit-learn==1.9.1 statsmodels==0.15.0 matplotlib==3.11.2 pillow==12.3.0
python3 -m pip install openpyxl plotly pypdf wordfreq playwright pymupdf
\end{verbatim}
On some Linux systems pip needs the flag \texttt{--break-system-packages} or a virtual environment (\texttt{python3 -m venv .venv} and then \texttt{source .venv/bin/activate}).

\subsection{JavaScript packages (npm)}
\begin{verbatim}
cd Software/v1.1.0/SWAFL_Desktop
npm install                      # electron ^31.0.0 and electron-builder ^24.13.3
npm run get-app                  # builds SWAFL.exe
npm install -g playwright        # for the test and screenshot scripts
npx playwright install chromium
\end{verbatim}
Leaflet, GeoTIFF.js and SheetJS are stored in the program folders and need no installation.

\subsection{Programs outside Python and JavaScript}
\begin{verbatim}
sudo apt-get install ffmpeg pandoc texlive-xetex texlive-latex-extra \
    texlive-fonts-extra texlive-bibtex-extra tesseract-ocr poppler-utils \
    fonts-dejavu fonts-texgyre
\end{verbatim}
On Windows the same programs come from their own installers (ffmpeg and pandoc and MiKTeX or TeX Live and Tesseract and poppler); Daisy 7.1.14 is already inside the program folder.

\subsection{Run the checks}
\begin{verbatim}
cd Manuscripts/reproducibility
bash run_all.sh ./work           # regenerates the numbers and figures
\end{verbatim}

\section{Packages inside the bundled libraries and inside Daisy}
A library that is copied into the program carries packages of its own. They are listed here so that the stack is complete from the program down to its smallest part. The libraries are minified single files, so the packages were identified from their code and their licence notices (not from a package list, which a minified file does not have); no version is claimed for them.
\subsection{GeoTIFF.js}
GeoTIFF.js (MIT licence, EOX IT Services) reads Sentinel-2 bands from remote GeoTIFF files by HTTP range requests. The program calls \texttt{GeoTIFF.fromUrl} (the maps and the monitoring of the field and the site chain) and \texttt{GeoTIFF.fromArrayBuffer}. Inside the file are these packages.
\begin{table}[H]\centering\small
\caption{Packages found inside \texttt{app/vendor/geotiff.js}.}\label{tab:geotiff}
\begin{tabular}{@{}>{\raggedright\arraybackslash}p{3.0cm}>{\raggedright\arraybackslash}p{3.0cm}>{\raggedright\arraybackslash}p{7.0cm}@{}}
\toprule Package & Evidence in the file & Role and why it is used in the build\\\midrule
pako (zlib in JavaScript) & Inflate and Z\_SYNC\_FLUSH code & decompresses DEFLATE tiles, a common compression of cloud-optimised files. Why: a tile cannot be read without it.\\
LERC (Esri, Apache 2.0) & the LERC decoder with its Esri notice & decodes the LERC compression of some rasters. Why: the same reader then works for more sources.\\
Zstandard decoder (zstddec) & \texttt{ZSTDDecoder} & decodes Zstandard tiles. Why: newer cloud-optimised files use it.\\
Half-float support (a float16 package) & \texttt{getFloat16} & reads 16-bit floating point samples. Why: some rasters store them.\\
Header parser (parse-headers style) & splitting of the header text at line ends & reads the \texttt{Content-Range} header of a range request. Why: the reader must know how long the whole file is.\\
Block cache (quick-lru style) & a maximum-size cache object & keeps recently read blocks. Why: fewer repeated requests for the same bytes.\\
XML helper (xml-utils style) & \texttt{findTagsByName} & reads the GDAL metadata of a file (scale and offset and no-data value). Why: the reflectance has to be scaled correctly.\\
Browser interfaces & \texttt{fetch} and \texttt{createImageBitmap} & range requests and the decoding of JPEG tiles by the browser. Why: no JPEG decoder has to be shipped.\\
\bottomrule\end{tabular}\end{table}
A package that GeoTIFF.js can use for worker threads (a web-worker shim) was searched for and not found, so the program reads the tiles in the main thread.

\subsection{SheetJS, Leaflet and Plotly.js}
\begin{itemize}[nosep,leftmargin=1.2em]
\item \textbf{SheetJS} (Apache 2.0) carries a table of character code pages (\texttt{cptable} 1.15.0, \texttt{DO\_NOT\_EXPORT\_CODEPAGE} set), which lets it read old Excel and CSV files in other encodings. Why: weather files from other countries do not use one encoding.
\item \textbf{Leaflet} 1.9.4 has no packages of its own; it is one file with its style sheet and images.
\item \textbf{Plotly.js} 4.1.1 (Results folder only) bundles its own graphics and mapping code, including the MapLibre GL code that its map charts use. It is not part of the program.
\end{itemize}

\subsection{Libraries shipped with Daisy}
The Daisy folder of the program (\texttt{daisy/bin}) carries the program and the libraries it was built with. They are listed from the file names of the folder.
\begin{table}[H]\centering\small
\caption{Files of the Daisy folder and what they are.}\label{tab:daisydll}
\begin{tabular}{@{}>{\raggedright\arraybackslash}p{4.2cm}>{\raggedright\arraybackslash}p{8.8cm}@{}}
\toprule File or package & Role and why it is used in the build\\\midrule
\texttt{daisy-bin.exe} & the Daisy program, started by the workbench. Why: the process model itself.\\
\texttt{libcore.dll} and \texttt{libwindows.dll} & the core and Windows parts of Daisy (judged from the file names). Why: Daisy was built as a program and libraries.\\
\texttt{libboost\_filesystem-mt.dll} (Boost) & file and folder handling. Why: Daisy reads and writes many files.\\
\texttt{libcxsparse.dll} and \texttt{libsuitesparseconfig.dll} (SuiteSparse) & sparse matrix solving. Why: the Richards and solute equations are solved implicitly with sparse matrices.\\
\texttt{libgomp-1.dll} (GNU OpenMP) & parallel loops. Why: speeds up the solvers on several cores.\\
\texttt{libstdc++-6.dll} and \texttt{libgcc\_s\_seh-1.dll} and \texttt{libwinpthread-1.dll} (GNU compiler runtime) & the C++ runtime. Why: the program was built with the GNU compilers for Windows.\\
\texttt{vcruntime140.dll} and \texttt{vcruntime140\_1.dll} (Microsoft) & the Visual C++ runtime. Why: needed by the Python runtime that Daisy carries.\\
\texttt{python313.dll} and \texttt{python.exe} and the \texttt{Lib} and \texttt{DLLs} folders (Python 3.13) & the Python runtime and its standard library. Why: Daisy can call Python functions inside a set-up.\\
\bottomrule\end{tabular}\end{table}

\section{Online services (optional)}
\begin{table}[H]\centering\small
\caption{Services the program can call. All are optional; without a connection every other number is computed on the computer.}\label{tab:svc}
\begin{tabular}{@{}>{\raggedright\arraybackslash}p{4.1cm}>{\raggedright\arraybackslash}p{9.0cm}@{}}
\toprule Service & Role and why it is used\\\midrule
Open-Meteo (forecast and archive and geocoding and elevation) & weather of any place (measured archive and forecast); place names; elevation. Why: free weather with no key, so the program needs no account.\\
Esri World Imagery and OpenStreetMap tiles & the satellite picture and the street map under the field. Why: open background maps so the user can see the field.\\
Microsoft Planetary Computer and Element 84 Earth Search (STAC) & Sentinel-2 L2A scenes for the vegetation and water indices. Why: open access to Sentinel-2 scenes in a standard catalogue (STAC).\\
SoilGrids (ISRIC) & soil of the field and its retention curves. Why: a soil estimate for any field in the world when the user has no soil data.\\
ESA WorldCover and Overpass (OpenStreetMap) & land cover and roads and water sources near the field. Why: tells the user what surrounds the field (roads and water) without a survey.\\
Nominatim and Photon & finding a place by name. Why: lets the user type a place name instead of coordinates.\\
\bottomrule\end{tabular}\end{table}

\section{Techniques used in the code}
\subsection{Structure of the program}
\begin{itemize}[nosep,leftmargin=1.2em]
\item \textbf{Two halves joined by a bridge.} The workbench (shell page) and the engine (inline frame) share state through one object that the bridge sends to the engine and reads back; a setting changed in Properties is applied to every layer, the dashboard, the report and the designs, so they cannot disagree.
\item \textbf{Canvas drawing.} Every picture (the field and the layers and the designs and the scientific figures) is drawn on a 2D canvas by code, so a figure can be exported as a picture or as numbers from the same data.
\item \textbf{Lazy loading and guards.} Heavy scripts load when a tab is first opened; a guard module stops a failed script from freezing the window; long runs are cut into steps so the window stays responsive.
\item \textbf{Provenance by construction.} Every number that is shown carries a label that leads to a numbered equation of the manual; the table of equations is generated from one file (\texttt{eqs.json}).
\item \textbf{The rule of the Daisy layers.} A Daisy layer shows a result only from a Daisy run: results are read back from the log files of the run, never typed in.
\item \textbf{Seeded randomness.} Projected seasons reuse the simulated seasons in a fixed random order from a seeded generator, so a projection is exactly repeatable.
\end{itemize}
\subsection{Methods of the models}
\begin{itemize}[nosep,leftmargin=1.2em]
\item Soil hydraulics from texture with the Saxton and Rawls pedotransfer functions; the van Genuchten and Mualem retention and conductivity; the Richards equation solved in the soil column tab of the workbench.
\item Surface irrigation: Kostiakov and Green and Ampt intake, Manning flow in the furrow, advance and recession by volume balance (Walker and Skogerboe).
\item Pressurised irrigation: Hazen and Williams friction (C of 140 and 150), speed limits, pump head with lift and friction and working head, emitter and sprinkler layouts.
\item Crop water: the FAO paper 56 daily balance with the crop coefficient method and the yield response of FAO paper 33; the FAO Penman and Monteith evapotranspiration.
\item Satellite chain: scene search by STAC, band reading by range requests, cloud mask from the scene classification, vegetation and water indices per pixel, zonal statistics inside the outlined field.
\item Experiment planning for Daisy: a catalogue of the factors Daisy can vary, a check of every design against it, an exact run budget, and a summary of the effect of every input from the logs.
\item Verification by independent references: worked examples and held-out simulations and stored Daisy output, with the defect found in the denitrification function reported in the articles.
\end{itemize}

\section{How the software is built}
\begin{table}[H]\centering\small
\caption{The build of the program.}\label{tab:build}
\begin{tabular}{@{}>{\raggedright\arraybackslash}p{3.7cm}>{\raggedright\arraybackslash}p{9.3cm}@{}}
\toprule Step & What happens and why\\\midrule
Start from source & \texttt{START\_SWAFL.bat} (or \texttt{.command}) starts the program from the folder; it finds Node.js or fetches a portable copy of about 30 to 40\,MB into \texttt{tools/node}. Why: a user can run the program before any build.\\
Install the build tools & \texttt{npm install} fetches Electron and electron-builder (24.13.3) listed in \texttt{package.json}. Why: the build tools stay out of the program and are needed only to make the file.\\
Build & \texttt{npm run get-app} runs electron-builder for Windows and writes \texttt{SWAFL.exe}. Why: one command makes the same file on every machine.\\
Targets & \texttt{portable} (one file and no installation) and \texttt{nsis} (an installer); a macOS \texttt{dmg} target is also defined. Why: a portable file for no-install use and an installer for those who want a shortcut.\\
Packed into the file & the files of the shell and \texttt{app/**/*} and the icon; as extra resources the manual PDF and the \texttt{daisy} folder. Why: the program must work offline with Daisy and the manual inside.\\
Web build & none needed: \texttt{npm run web} serves the same \texttt{app} folder with \texttt{serve.js}. Why: the same files serve the browser, so the web and desktop versions cannot drift apart.\\
\bottomrule\end{tabular}\end{table}
The program is not code-signed, so Windows may show a SmartScreen notice, which is cleared with More info and then Run anyway.

\section{How the documents are built}
\begin{itemize}[nosep,leftmargin=1.2em]
\item \textbf{The manual.} The chapters are Markdown files. A Python tool (\texttt{build\_single.py}) joins them, adds the icons of the layers and the equation table, and \texttt{pandoc} 3.1.3 with the XeLaTeX engine (TeX Live 2023) makes one PDF with DejaVu Serif and DejaVu Sans Mono. Other tools write the index of sections that the program uses to open the manual at the page in use, and a short script with PyMuPDF (run at build time and not stored in the repository) reduces the size of the pictures so that the PDF opens from GitHub on any device.
\item \textbf{The screenshots of the manual.} Taken from the running program with Playwright and headless Chromium, then annotated by Python scripts (\texttt{annot.py}, \texttt{new\_capture.py}).
\item \textbf{The articles and supplements.} \LaTeX\ with XeLaTeX and BibTeX (\texttt{build.sh}); the tables and the numbers are written into the source from the data files by scripts, so that the text, the tables and the figures use the same numbers.
\item \textbf{The figures of the series.} Python in the scripts of the repository, with the packages of the Python table above.
\end{itemize}

\section{Tools used to test and check}
\begin{itemize}[nosep,leftmargin=1.2em]
\item Playwright with headless Chromium drives the real program: every layer is opened and run, the dashboard and the figures are read back, and page errors in the console are counted (the target is none).
\item Node.js scripts run the equations of the designs outside the page and compare them with worked examples.
\item Python scripts compare results with the stored Daisy output (\texttt{compare\_results.py}, \texttt{audit\_numbers.py}) and check the data against a table of checksums (\texttt{CHECKSUMS.sha256}).
\item A local static server (\texttt{python -m http.server}) serves the program for these tests so that the same files are tested as are shipped.
\end{itemize}

\section{Licences and credit}
The code of SWAFL and the articles are by Naziru Halilu and are released under the MIT licence of the repository. Leaflet and GeoTIFF.js and SheetJS keep their own licences and notices in the folders \texttt{app/vendor} and \texttt{app/engine/vendor}. Daisy keeps its licence in its own folder and is credited to the Daisy project of the University of Copenhagen \citep{Hansen2012}. The online services keep the terms of their owners and are credited where their data appear.

\section{Where to find it}
All parts are in the repository \emph{Coupled Water Fluxes, Pore-Scale Dynamics, Precision Irrigation, Mulching Feedback Loops}: the program in \texttt{Software/v1.1.0/SWAFL\_Desktop}, the manual in \texttt{Software/v1.1.0/Manual}, the articles and supplements in \texttt{Manuscripts/pdf}, and the data in \texttt{Data} and \texttt{Results}. The way to rebuild the program is given in the install chapter of the manual.

\bibliographystyle{plainnathl}
\bibliography{refs}
\end{document}
